% =====================================================================
%  Peptide Identification by Exclusion Search:
%  A Database Algorithm Whose Output Is a Region, Whose Termination
%  Is Closure, and Whose Refusal Is a Verdict
%
%  Self-contained. Every construction used is defined here.
% =====================================================================
\documentclass[11pt,a4paper]{article}

\usepackage[utf8]{inputenc}
\usepackage[T1]{fontenc}
\usepackage{amsmath,amssymb,amsfonts,amsthm}
\usepackage{mathtools}
\usepackage[margin=1in]{geometry}
\usepackage{graphicx}
\usepackage{booktabs}
\usepackage{array}
\usepackage{enumitem}
\usepackage{lmodern}
\usepackage[protrusion=true,expansion=false]{microtype}
\usepackage[numbers,sort&compress]{natbib}
\usepackage[colorlinks=true,linkcolor=blue,citecolor=blue,urlcolor=blue]{hyperref}
\usepackage[capitalise,nameinlink]{cleveref}
\usepackage{algorithm}
\usepackage{algpseudocode}
\usepackage{siunitx}

% ---------------------------------------------------------------------
%  Theorem environments
% ---------------------------------------------------------------------
\newtheorem{theorem}{Theorem}[section]
\newtheorem{lemma}[theorem]{Lemma}
\newtheorem{corollary}[theorem]{Corollary}
\newtheorem{proposition}[theorem]{Proposition}
\theoremstyle{definition}
\newtheorem{definition}[theorem]{Definition}
\newtheorem{axiom}[theorem]{Axiom}
\newtheorem{construction}[theorem]{Construction}
\newtheorem{example}[theorem]{Example}
\theoremstyle{remark}
\newtheorem{remark}[theorem]{Remark}

% ---------------------------------------------------------------------
%  Notation
% ---------------------------------------------------------------------
\newcommand{\Gph}{\mathcal{G}}
\newcommand{\wt}{w}
\newcommand{\med}{\mathfrak{m}}
\newcommand{\flo}{\beta}
\newcommand{\cut}{\operatorname{cut}}
\newcommand{\sep}{\sigma}
\newcommand{\dep}{\delta}
\newcommand{\key}{\kappa}
\newcommand{\resid}{R}
\newcommand{\clos}{\mathrm{cl}}
\newcommand{\amb}{\mathcal{A}}
\newcommand{\Cat}{\mathrm{cell}}
\newcommand{\dem}{D}
\newcommand{\Sk}{S_k}
\newcommand{\St}{S_t}
\newcommand{\Se}{S_e}
\newcommand{\Scoord}{\mathbf{S}}
\newcommand{\Prot}{\mathcal{P}}
\newcommand{\Pep}{\mathcal{Q}}
\newcommand{\Spec}{\mathcal{X}}
\newcommand{\DB}{\mathcal{B}}
\newcommand{\pot}{\kappa}
\DeclareMathOperator{\prom}{prom}
\DeclareMathOperator*{\argmin}{arg\,min}

\title{\textbf{Peptide Identification by Exclusion Search:\\[2pt]
A Database Algorithm Whose Output Is a Region,\\
Whose Termination Is Closure,\\
and Whose Refusal Is a Verdict}}

\author{Kundai Farai Sachikonye\\
\texttt{kundai.sachikonye@wzw.tum.de}}

\date{\today}

\begin{document}
\maketitle

% =====================================================================
\begin{abstract}
\noindent
We give a complete, implementable database search algorithm for peptide
tandem mass spectrometry that differs from existing engines in three
structural respects. Its output is a \emph{region}---a set of candidate
sequences indistinguishable under committed evidence---rather than a
ranked list filtered by a confidence threshold. Its termination
condition is \emph{closure}---the state in which no further evidence
source available to the search can move the region---rather than a score
exceeding a threshold. And it has two terminal states, closure and
\emph{honest decline}, with no third state in which the algorithm is
still running and has no internal signal of its own futility.

The algorithm constructs, for each query spectrum, a contact graph whose
items are candidate peptides surviving a precursor mass filter and whose
edges are committed by a chain of five catalysts: precursor mass
exclusion, enzyme specificity exclusion, fragment ion matching,
S-Entropy coordinate similarity, and ladder topology matching. Each
catalyst has a computable catalytic power, and their combination follows
a multiplicative composition law with a precise diminishing-returns
profile. The claim region after all catalysts is the identity of the
peptide-spectrum match: it may contain one candidate (unique
identification), more than one (a terminal region such as $\{L,I\}$
variants), or the algorithm may decline (insufficient evidence). In all
three cases the output is a verdict, not a score.

We prove that the algorithm terminates in at most $D_0$ catalyst
applications, where $D_0$ is the initial demand; that closure is
strictly stronger than any FDR threshold; that an isobaric pair is a
correct closed answer under mass evidence; and that agreement between an
independently constructed instrument graph and the database mapping
graph is a relabelling invariant not obtainable from either alone.

The algorithm runs in $O(|\Pep_M| \cdot L \cdot m)$ time per spectrum,
where $|\Pep_M|$ is the number of precursor-mass-filtered candidates,
$L$ is maximal peptide length, and $m$ is the number of observed peaks.
It requires no training data, no decoy database, and no probability
model. Where it cannot identify, it says so and says why.
\end{abstract}

\noindent\textbf{Keywords:} peptide identification; database search;
exclusion; closure; contact graph; S-Entropy; tandem mass spectrometry;
region-valued identity; honest decline.

\tableofcontents
\newpage

% =====================================================================
\section{Introduction}
\label{sec:intro}
% =====================================================================

\subsection{The problem}

Peptide identification by tandem mass spectrometry proceeds, in every
widely used engine, by the same logic: generate theoretical spectra from
a protein database, score each against the observed spectrum, rank by
score, and filter by a false discovery rate (FDR) threshold computed
from a decoy
database~\citep{eng1994approach,perkins1999probability,cox2008maxquant,%
elias2007target,kall2007semi,the2016fast}. The output is a ranked list
of peptide-spectrum matches (PSMs), each carrying a score and a
posterior error probability.

This logic has three structural properties that are independent of any
particular scoring function.

\paragraph{The output is a point.} Each PSM is a single sequence
assignment. Where two sequences are indistinguishable (leucine/isoleucine,
isobaric modifications at different sites), the engine either picks one
arbitrarily or reports the ambiguity as a quality annotation, not as the
identity itself.

\paragraph{Termination is by threshold.} The engine stops reporting when
the score falls below a threshold or the FDR exceeds a cutoff. The
threshold is external: it is not computed from the evidence that bears on
any individual spectrum, and different thresholds yield different
identifications from the same data.

\paragraph{Silence is the only refusal.} A spectrum that matches nothing
above threshold simply does not appear in the output. The engine does
not say \emph{why} it failed: whether the peptide is absent from the
database, whether the spectrum quality is insufficient, whether the
fragmentation pattern is inconsistent with any candidate, or whether two
candidates tie. All four situations produce the same output: nothing.

We present an algorithm in which all three properties are replaced.

\subsection{What this algorithm does differently}

\begin{enumerate}[label=(\roman*),leftmargin=2.2em,itemsep=2pt]
  \item The output is a \emph{region}: a set of candidate sequences
  that are indistinguishable under all evidence the algorithm has
  committed. The set may contain one element (unique identification),
  more than one (a terminal region), or the algorithm may decline.
  \item Termination is by \emph{closure}: the state in which no
  further catalyst available to the search can move the region. Closure
  is strictly stronger than any confidence threshold
  (\Cref{thm:closure-stronger}).
  \item There are exactly two terminal states: closure and honest
  decline. Decline carries a structured verdict stating which evidence
  source was responsible (\Cref{thm:dichotomy}).
\end{enumerate}

\subsection{Relation to existing work}

The algorithm's scoring components---fragment ion matching, precursor
mass filtering, enzyme specificity---are standard and individually
well understood~\citep{eng1994approach,paizs2005fragmentation,%
olsen2004trypsin,krokhin2006sequence}. The S-Entropy coordinate space
is developed in~\citet{sachikonye2025molecular}; the phase-lock
network topology in~\citet{sachikonye2025fragmentation}. The
theoretical framework of exclusion, closure, and region-valued identity
is developed in~\citet{sachikonye2025peptide}. What is new here is the
synthesis: an implementable algorithm that threads these components
through the exclusion machinery, with proofs of termination, closure
strength, and complexity.

\subsection{What this algorithm does not do}

It does not replace a scoring function with a better one. It replaces
the \emph{kind of object} the search returns (a region instead of a
point) and the \emph{kind of condition} under which it stops (closure
instead of a threshold). It is therefore orthogonal to improvements in
peak matching, intensity prediction, or machine-learned
scoring~\citep{tran2017novo,zhou2017pdeep}: any of those may serve as a
catalyst within this framework, and the closure and region guarantees
hold regardless of which catalysts are used.

It does not eliminate the database. It uses the database as one evidence
source among several, and the exclusion framework specifies what the
database contributes (a mapping graph) and what it does not (an
instrument graph).

% =====================================================================
\section{Preliminaries}
\label{sec:prelim}
% =====================================================================

We state the definitions and results from the exclusion framework that
the algorithm requires. Proofs are
in~\citet{sachikonye2025peptide}; we reproduce only what is needed
for self-containedness.

\begin{definition}[Contact graph]
\label{def:contact}
A \emph{contact graph} is a finite weighted graph
$\Gph=(V,E,\wt)$ with $\wt:E\to\mathbb{R}_{>0}$, together with a
distinguished vertex $\med\in V$ called the \emph{medium}, adjacent to
every other vertex. Vertices in $V\setminus\{\med\}$ are \emph{items}.
\end{definition}

\begin{definition}[Separation cost, cut key]
\label{def:sep}
For item $v$, the \emph{separation cost} $\sep(v)$ is the minimum
$v$-$\med$ cut weight. The \emph{depth} $\dep(v)=|S^{\ast}(v)|$ is
the minimising side's cardinality. The \emph{cut key}
$\key(v)=(\sep(v),\dep(v))$ is the relabelling-invariant identity
of~$v$.
\end{definition}

\begin{theorem}[Floor]
\label{thm:floor}
There exists $\flo>0$ such that $\sep(v)\ge\flo$ for every item $v$.
\end{theorem}

\begin{definition}[Catalyst]
\label{def:catalyst}
A \emph{catalyst} $c$ commits one or more edges into a contact graph.
It has \emph{catalytic power}
\[
  \pot(c,v,x^{\ast})
  \;=\;
  \frac{a(v,x^{\ast})-a(c(v),x^{\ast})}{a(v,x^{\ast})-\flo/\Omega}
  \;\in\;[0,1],
\]
where $a(\cdot,\cdot)$ is the normalised alignment functional.
\end{definition}

\begin{theorem}[Multiplicative composition]
\label{thm:mult}
For catalysts $c_1,\dots,c_n$ applied sequentially with powers
$\pot_1,\dots,\pot_n$, the composite power is
$1-\prod_{i=1}^{n}(1-\pot_i)$.
\end{theorem}

\begin{definition}[Claim region, closure]
\label{def:closure}
The \emph{claim region} $\amb(v)$ of item $v$ is the set of candidates
whose cut key agrees with $\key(v)$ within tolerance. A determination
is \emph{closed} if every available catalyst leaves $\amb(v)$ unchanged.
\end{definition}

\begin{theorem}[Closure is strictly stronger than a threshold]
\label{thm:closure-stronger}
For any confidence functional $\mathrm{conf}$ and threshold $\tau$:
(i) $\clos(v)$ does not imply $\mathrm{conf}(v)>\tau$;
(ii) $\mathrm{conf}(v)>\tau$ does not imply $\clos(v)$;
(iii) $\clos(v)$ is decidable from the graph; $\tau$ is not.
\end{theorem}

\begin{theorem}[Termination dichotomy]
\label{thm:dichotomy}
A determination has exactly two outcomes: closure (in at most $D_0$
catalyst applications) or honest decline. There is no third outcome.
\end{theorem}

\begin{theorem}[Two-graph agreement]
\label{thm:agreement}
If an instrument graph $\Gph_I$ and a mapping graph $\Gph_M$ both close
on item $v$, their claim-region intersection is a relabelling invariant
not obtainable from either alone.
\end{theorem}

% =====================================================================
\section{The S-Entropy Coordinate Space}
\label{sec:sentropy}
% =====================================================================

\begin{definition}[Amino acid coordinate map]
\label{def:sentropy}
Define $\phi_{\mathrm{AA}}:\mathcal{R}\to[0,1]^3$ mapping each amino
acid residue $r\in\mathcal{R}$ to coordinates $(\Sk,\St,\Se)$ by:
\begin{align}
  \Sk(r) &= \frac{H(r)-H_{\min}}{H_{\max}-H_{\min}}, \label{eq:sk}\\[4pt]
  \St(r) &= \frac{V_{\mathrm{vdW}}(r)-V_{\min}}{V_{\max}-V_{\min}},
  \label{eq:st}\\[4pt]
  \Se(r) &= \frac{|\chi(r)|-|\chi|_{\min}}{|\chi|_{\max}-|\chi|_{\min}},
  \label{eq:se}
\end{align}
where $H(r)$ is the Kyte--Doolittle
hydrophobicity~\citep{kyte1982simple}, $V_{\mathrm{vdW}}(r)$ the van
der Waals volume~\citep{zamyatnin1972protein}, and $\chi(r)$ the
absolute net charge at pH~7.0. Extrema are over the 20 standard amino
acids.
\end{definition}

\begin{definition}[Peptide S-Entropy profile]
\label{def:profile}
For a peptide $p=r_1 r_2\cdots r_L$, the \emph{S-Entropy profile} is
the sequence $(\phi_{\mathrm{AA}}(r_1),\dots,\phi_{\mathrm{AA}}(r_L))
\in([0,1]^3)^L$, and the \emph{summary coordinate} is
\[
  \bar{\Scoord}(p)
  \;=\;
  \frac{1}{L}\sum_{i=1}^{L}\phi_{\mathrm{AA}}(r_i).
\]
\end{definition}

\begin{table}[t]
\centering
\small
\caption{S-Entropy coordinates for the 20 standard amino acids, computed
from \Cref{def:sentropy}. Values rounded to three decimal places.}
\label{tab:aa-coords}
\begin{tabular}{@{}lccclccc@{}}
\toprule
\textbf{AA} & $\Sk$ & $\St$ & $\Se$ &
\textbf{AA} & $\Sk$ & $\St$ & $\Se$ \\
\midrule
G & 0.456 & 0.000 & 0.000 & Y & 0.322 & 0.812 & 0.000 \\
A & 0.578 & 0.107 & 0.000 & W & 0.389 & 1.000 & 0.000 \\
V & 0.822 & 0.393 & 0.000 & D & 0.211 & 0.339 & 0.500 \\
L & 0.867 & 0.536 & 0.000 & E & 0.211 & 0.464 & 0.500 \\
I & 0.900 & 0.536 & 0.000 & N & 0.211 & 0.375 & 0.000 \\
P & 0.356 & 0.321 & 0.000 & Q & 0.211 & 0.500 & 0.000 \\
F & 0.778 & 0.750 & 0.000 & K & 0.178 & 0.571 & 0.500 \\
M & 0.567 & 0.536 & 0.000 & R & 0.000 & 0.714 & 1.000 \\
S & 0.344 & 0.143 & 0.000 & H & 0.178 & 0.571 & 0.275 \\
T & 0.344 & 0.268 & 0.000 & C & 0.711 & 0.214 & 0.000 \\
\bottomrule
\end{tabular}
\end{table}

\begin{remark}[Leucine and isoleucine]
\label{rem:li}
L and I share $\St=0.536$ and $\Se=0.000$ but differ at $\Sk$:
$0.867$ vs.\ $0.900$. The separation is small but nonzero, and the
coordinate space records it. Whether it is resolvable depends on the
catalysts available; \Cref{prop:isobaric} addresses this precisely.
\end{remark}

% =====================================================================
\section{The Ladder Topology Features}
\label{sec:topology}
% =====================================================================

We extract from each spectrum a set of features that are invariant
across instrument platforms, following~\citet{sachikonye2025fragmentation}.

\begin{definition}[Ladder topology]
\label{def:topology}
Given an MS/MS spectrum $\Spec$ with precursor mass $M$ and charge $z$,
and a candidate peptide $p$ of length $L$, the \emph{ladder topology}
is the quadruple
$T(\Spec,p)=(T_b,T_y,T_c,T_r)$ defined by:
\begin{align}
  T_b &= \frac{|\{i : b_i^{\mathrm{theo}} \text{ matched in } \Spec\}|}{L-1},
  \label{eq:tb}\\[2pt]
  T_y &= \frac{|\{i : y_i^{\mathrm{theo}} \text{ matched in } \Spec\}|}{L-1},
  \label{eq:ty}\\[2pt]
  T_c &= \frac{|\{i : b_i \text{ and } y_{L-i} \text{ both matched}\}|}{L-1},
  \label{eq:tc}\\[2pt]
  T_r &= 1 - \frac{\mathrm{Var}(\Delta m_{b})}{(\mathrm{Mean}(\Delta m_{b}))^2},
  \label{eq:tr}
\end{align}
where $b_i^{\mathrm{theo}}=\sum_{j=1}^{i}m(r_j)+1.008$ and
$y_i^{\mathrm{theo}}=M-b_{L-i}^{\mathrm{theo}}+19.018$
are the theoretical b- and y-ion masses, a match is declared when
$|m_{\mathrm{obs}}-m_{\mathrm{theo}}|<\epsilon_f$ for fragment
tolerance $\epsilon_f$, and $\Delta m_b$ is the vector of consecutive
b-ion spacings among matched ions.
\end{definition}

\begin{remark}
$T_b$ and $T_y$ measure coverage, $T_c$ measures complementarity,
and $T_r$ measures regularity of the ladder. All four lie in $[0,1]$
and are dimensionless. The claim
in~\citet{sachikonye2025fragmentation} is that these achieve
CV $<2.1\%$ across four instrument platforms; we use them here as
catalyst inputs without reproducing the platform-independence validation.
\end{remark}

% =====================================================================
\section{Database Preparation}
\label{sec:database}
% =====================================================================

\begin{construction}[In-silico digestion and index]
\label{con:digest}
Given a protein database $\DB$ (a finite set of amino acid sequences)
and digestion parameters (enzyme $\mathcal{E}$, maximum missed
cleavages $c_{\max}$, minimum/maximum peptide length $[L_{\min},L_{\max}]$):
\begin{enumerate}[leftmargin=2.0em,itemsep=1pt]
  \item Cleave each protein at enzyme-specific sites, retaining
  peptides of length in $[L_{\min},L_{\max}]$ with at most
  $c_{\max}$ internal cleavage sites.
  \item For each peptide $p$, compute and store:
  \begin{itemize}[itemsep=1pt]
    \item the neutral monoisotopic mass $M(p)$;
    \item the S-Entropy summary coordinate $\bar{\Scoord}(p)$;
    \item the full S-Entropy profile
    $(\phi_{\mathrm{AA}}(r_1),\dots,\phi_{\mathrm{AA}}(r_L))$;
    \item the theoretical b- and y-ion mass arrays;
    \item the mapping set $\amb(p)=\{P\in\DB : p\text{ occurs in }P\}$.
  \end{itemize}
  \item Build a sorted index on $M(p)$ for binary-search precursor
  filtering.
\end{enumerate}
\end{construction}

\begin{remark}[Fixed modifications]
Cysteine carbamidomethylation ($+57.021$~Da) or other fixed
modifications are applied at digestion time: $M(p)$ includes
them, and $\Sk$, $\St$, $\Se$ are computed on the modified residue.
Variable modifications are handled as additional candidates in the
digest, each generating a separate entry with modified mass and
coordinate.
\end{remark}

% =====================================================================
\section{The Exclusion Search Algorithm}
\label{sec:algorithm}
% =====================================================================

\subsection{Overview}

The algorithm takes as input a centroided MS/MS spectrum
$\Spec=\{(m_j,I_j)\}_{j=1}^{m}$, precursor mass $M$, charge state $z$,
and the indexed database of \Cref{con:digest}. It returns a verdict:
either a closed claim region $\amb(v)$ (which may have $|\amb|=1$ or
$|\amb|>1$), or an honest decline with a diagnostic payload.

The search proceeds through a chain of five catalysts, each committing
edges in a contact graph constructed per spectrum. The chain is ordered
by decreasing exclusion breadth: cheap, broad filters first; expensive,
narrow discriminators last.

\subsection{Catalyst 1: Precursor mass exclusion}

\begin{definition}[Precursor filter]
\label{def:prec}
Given precursor tolerance $\epsilon_p$ (in Da or ppm), define:
\[
  \Pep_M \;=\; \{p\in\Pep : |M(p)-M|<\epsilon_p\}.
\]
All peptides $p\notin\Pep_M$ are excluded. The exclusion weight of this
catalyst is $|\Pep|-|\Pep_M|$.
\end{definition}

\begin{remark}
$\Pep_M$ is computed by binary search on the sorted mass index in
$O(\log|\Pep|+|\Pep_M|)$ time. For $\epsilon_p=10$~ppm on a
typical human proteome digest ($|\Pep|\approx 10^6$), $|\Pep_M|$ is
typically $10^1$--$10^3$.
\end{remark}

\subsection{Catalyst 2: Enzyme specificity exclusion}

\begin{definition}[Enzyme filter]
\label{def:enzyme}
For enzyme $\mathcal{E}$ with cleavage specificity (e.g.\ trypsin:
C-terminal to K/R, not before P), exclude candidates violating the
specificity at either terminus. For semi-specific searches, require
specificity at one terminus only.
\end{definition}

This catalyst is applied during digestion (\Cref{con:digest}) and costs
no additional runtime. Its exclusion weight is the number of peptides
removed by the specificity constraint.

\subsection{Catalyst 3: Fragment ion matching}

This is the primary discriminator and the most expensive per candidate.

\begin{definition}[Fragment match score]
\label{def:frag}
For candidate $p$ of length $L$ against spectrum $\Spec$, define:
\begin{align}
  n_b(p) &= |\{i\in[1,L-1] : \exists\,j,\;
  |m_j-b_i^{\mathrm{theo}}|<\epsilon_f\}|,\\[2pt]
  n_y(p) &= |\{i\in[1,L-1] : \exists\,j,\;
  |m_j-y_i^{\mathrm{theo}}|<\epsilon_f\}|,\\[2pt]
  n_{\mathrm{total}}(p) &= n_b(p)+n_y(p).
\end{align}
The \emph{fragment coverage} is $f(p)=n_{\mathrm{total}}(p)/(2(L-1))$.
\end{definition}

\begin{definition}[Fragment catalyst]
\label{def:frag-catalyst}
The fragment catalyst commits, for each candidate $p\in\Pep_M$, an edge
of weight $f(p)$ between $p$ and the medium $\med$. Its catalytic power
toward the target (the true peptide generating the spectrum) is
monotone in $f$: higher coverage gives higher power.
\end{definition}

\begin{remark}[Intensity is optional]
\Cref{def:frag} counts matched ions and does not weight by intensity.
Intensity-weighted variants are compatible with the framework---they
change the edge weight function, not the closure logic---but we use the
unweighted count as the base catalyst for platform independence
(\Cref{sec:topology}).
\end{remark}

\subsection{Catalyst 4: S-Entropy coordinate similarity}

\begin{definition}[Coordinate catalyst]
\label{def:coord-catalyst}
For candidate $p$ and the spectrum's \emph{empirical} S-Entropy
profile reconstructed from matched fragment masses, define:
\[
  d_S(p) \;=\; \|\bar{\Scoord}(p)-\bar{\Scoord}_{\mathrm{emp}}\|_2,
\]
where $\bar{\Scoord}_{\mathrm{emp}}$ is computed from the amino acid
residues implied by consecutive matched b-ion mass differences. The
catalyst commits an edge of weight $\max(0,\,1-d_S(p)/d_{\max})$
between $p$ and $\med$, where $d_{\max}=\sqrt{3}$ is the diagonal of
the unit cube.
\end{definition}

\begin{remark}
The empirical profile is available only where consecutive b- or y-ions
are matched, so it is partial. Where no consecutive pair is matched,
this catalyst is inert ($\pot=0$) and contributes nothing, which is
the correct behaviour: an inert catalyst does not claim evidence it
does not have.
\end{remark}

\subsection{Catalyst 5: Ladder topology matching}

\begin{definition}[Topology catalyst]
\label{def:topo-catalyst}
For candidate $p$ against spectrum $\Spec$, compute the topology
quadruple $T(\Spec,p)=(T_b,T_y,T_c,T_r)$ of \Cref{def:topology}.
The topology score is:
\[
  t(p) \;=\; \frac{1}{4}(T_b+T_y+T_c+T_r).
\]
The catalyst commits an edge of weight $t(p)$ between $p$ and $\med$.
\end{definition}

\begin{remark}[Independence from Catalyst~3]
Catalyst~3 counts matched ions; Catalyst~5 evaluates the \emph{pattern}
of matches (complementarity, regularity). Two candidates with identical
fragment counts may have different topology scores if one's matches
are clustered and the other's are evenly spaced. The two catalysts are
therefore not redundant, and their combination is governed by the
multiplicative law (\Cref{thm:mult}).
\end{remark}

\subsection{Contact graph construction}

\begin{construction}[Per-spectrum contact graph]
\label{con:graph}
For spectrum $\Spec$ with $|\Pep_M|=n$ candidates surviving precursor
filtering:
\begin{enumerate}[leftmargin=2.0em,itemsep=1pt]
  \item Create vertex set $V=\{p_1,\dots,p_n,\med\}$.
  \item For each $p_i$, connect $p_i$ to $\med$ with weight
  $w_0=\flo$ (the medium-incident floor edge).
  \item Apply catalysts 3, 4, 5 in sequence. Each catalyst commits
  additional edges between items (inter-candidate edges of weight
  equal to their pairwise S-Entropy distance) and modifies the
  item-medium edge weights (by adding the catalyst's score to the
  existing weight). Formally, after all catalysts:
  \[
    \wt(\{p_i,\med\}) \;=\; \flo + f(p_i) + (1-d_S(p_i)/d_{\max})
    + t(p_i).
  \]
  \item For each pair $p_i,p_j$ of candidates, commit an inter-item
  edge of weight $\|\bar{\Scoord}(p_i)-\bar{\Scoord}(p_j)\|_2$ if
  this exceeds $\flo$.
\end{enumerate}
\end{construction}

\subsection{Claim region computation}

\begin{construction}[Claim region]
\label{con:claim}
After committing all catalyst edges:
\begin{enumerate}[leftmargin=2.0em,itemsep=1pt]
  \item For each item $p_i$, compute $\sep(p_i)$ as the minimum
  $p_i$-$\med$ cut weight, by maximum
  flow~\citep{fordfulkerson1956,goldberg1988new}.
  \item Compute the cut key $\key(p_i)=(\sep(p_i),\dep(p_i))$.
  \item Let $p^{\ast}=\argmin_{p_i}\sep(p_i)$. The claim region is:
  \[
    \amb(\Spec) \;=\; \{p_i : \key(p_i)=\key(p^{\ast})\}.
  \]
\end{enumerate}
\end{construction}

\begin{remark}
$\amb(\Spec)$ is the set of candidates that are indistinguishable from
the best-separated item under the committed evidence. If
$|\amb(\Spec)|=1$, the identification is unique. If $|\amb(\Spec)|>1$,
the set is the correct answer: it names exactly the candidates the
evidence cannot tell apart.
\end{remark}

\subsection{Closure check}

\begin{construction}[Closure]
\label{con:closure}
After computing $\amb(\Spec)$, check closure:
\begin{enumerate}[leftmargin=2.0em,itemsep=1pt]
  \item For each catalyst $c$ in $\{3,4,5\}$, simulate the effect of
  \emph{re-applying} $c$ with perturbed parameters (e.g.\ tighter
  fragment tolerance, additional charge states for b/y ions).
  \item If no perturbation changes $\amb(\Spec)$, declare
  $\clos(\Spec)$: the determination is closed.
  \item If a perturbation changes $\amb(\Spec)$ and a catalyst
  implementing that perturbation is available (e.g.\ MS$^3$ data,
  ion mobility), apply it and recompute. Otherwise, proceed to
  decline.
\end{enumerate}
\end{construction}

\subsection{Decline and verdict}

\begin{definition}[Verdict]
\label{def:verdict}
The output of the algorithm on spectrum $\Spec$ is a \emph{verdict}
$\mathcal{V}(\Spec)$, one of:
\begin{center}
\begin{tabular}{@{}lp{0.65\textwidth}@{}}
\toprule
\textbf{Verdict} & \textbf{Payload} \\
\midrule
\textsc{identified} & $\amb(\Spec)$ (a set, $|\amb|\ge1$), closure
flag, composite catalytic power \\
\textsc{isobaric} & $\amb(\Spec)$ with $|\amb|>1$ and all members
differing only at L/I positions \\
\textsc{declined-no-candidate} & $\Pep_M=\emptyset$: no peptide in
the database matches the precursor mass \\
\textsc{declined-no-evidence} & $\Pep_M\neq\emptyset$ but all
catalysts are inert ($\pot=0$): e.g.\ no fragment ions matched \\
\textsc{declined-unclosed} & $\amb(\Spec)$ exists but a known catalyst
would move it and is unavailable \\
\bottomrule
\end{tabular}
\end{center}
\end{definition}

\begin{proposition}[Exhaustiveness]
\label{prop:exhaust}
Every spectrum receives exactly one verdict. The five cases are
mutually exclusive and jointly exhaustive.
\end{proposition}

\begin{proof}
The cases are ordered by evaluation: (1) $\Pep_M=\emptyset$ is checked
first; (2) all catalysts inert is checked next; (3) closure is checked;
(4) if closed with $|\amb|>1$ at L/I-only positions, verdict is
\textsc{isobaric}; otherwise (5) \textsc{identified} if closed,
\textsc{declined-unclosed} if not.
\end{proof}

\subsection{The full algorithm}

\begin{algorithm}[H]
\caption{Exclusion Search}
\label{alg:main}
\begin{algorithmic}[1]
\Require Spectrum $\Spec$, precursor mass $M$, charge $z$,
         indexed database $(\Pep,\mathrm{index})$,
         tolerances $\epsilon_p,\epsilon_f$, floor $\flo$
\Ensure Verdict $\mathcal{V}(\Spec)$
\Statex
\State $\Pep_M \gets \{p\in\Pep : |M(p)-M|<\epsilon_p\}$
  \Comment{Catalyst 1: precursor filter}
\If{$\Pep_M = \emptyset$}
  \State \Return \textsc{declined-no-candidate}
\EndIf
\Statex
\State \Comment{--- Build contact graph ---}
\State $V \gets \Pep_M \cup \{\med\}$
\For{each $p_i \in \Pep_M$}
  \State $f_i \gets \textsc{FragmentMatch}(p_i, \Spec, \epsilon_f)$
    \Comment{Catalyst 3}
  \State $d_i \gets \textsc{CoordDist}(p_i, \Spec)$
    \Comment{Catalyst 4}
  \State $t_i \gets \textsc{TopoScore}(p_i, \Spec, \epsilon_f)$
    \Comment{Catalyst 5}
  \State $\wt(\{p_i,\med\}) \gets \flo + f_i + \max(0,1-d_i/\sqrt{3})
    + t_i$
\EndFor
\Statex
\If{$\max_i(f_i + t_i) = 0$}
  \State \Return \textsc{declined-no-evidence}
\EndIf
\Statex
\State \Comment{--- Commit inter-item edges ---}
\For{each pair $(p_i, p_j)$ with $i<j$}
  \State $w_{ij} \gets \|\bar{\Scoord}(p_i)-\bar{\Scoord}(p_j)\|_2$
  \If{$w_{ij} > \flo$}
    \State commit edge $\{p_i,p_j\}$ with weight $w_{ij}$
  \EndIf
\EndFor
\Statex
\State \Comment{--- Compute claim region ---}
\For{each $p_i \in \Pep_M$}
  \State $\sep(p_i) \gets \textsc{MaxFlow}(p_i, \med, \Gph)$
  \State $\dep(p_i) \gets |S^{\ast}(p_i)|$
  \State $\key(p_i) \gets (\sep(p_i), \dep(p_i))$
\EndFor
\State $p^{\ast} \gets \argmin_{p_i} \sep(p_i)$
\State $\amb(\Spec) \gets \{p_i : \key(p_i)=\key(p^{\ast})\}$
\Statex
\State \Comment{--- Closure check ---}
\State $\mathrm{closed} \gets \textsc{CheckClosure}(\amb(\Spec),
  \Gph, \text{available catalysts})$
\If{$\mathrm{closed}$ \textbf{and}
    \textsc{IsLIOnly}($\amb(\Spec)$)}
  \State \Return \textsc{isobaric}($\amb(\Spec)$)
\ElsIf{$\mathrm{closed}$}
  \State \Return \textsc{identified}($\amb(\Spec)$,
    composite power)
\Else
  \State \Return \textsc{declined-unclosed}($\amb(\Spec)$,
    blocking catalyst)
\EndIf
\end{algorithmic}
\end{algorithm}

\begin{algorithm}[H]
\caption{CheckClosure}
\label{alg:closure}
\begin{algorithmic}[1]
\Require Claim region $\amb$, contact graph $\Gph$,
         list of available catalysts $C$
\Ensure Boolean
\For{each catalyst $c \in C$ not yet applied}
  \State $\Gph' \gets$ apply $c$ to $\Gph$
  \State $\amb' \gets$ recompute claim region on $\Gph'$
  \If{$\amb' \neq \amb$}
    \State \Return \textbf{false}
  \EndIf
\EndFor
\State \Return \textbf{true}
\end{algorithmic}
\end{algorithm}

% =====================================================================
\section{The Two-Graph Construction}
\label{sec:twograph}
% =====================================================================

The five catalysts of \Cref{sec:algorithm} produce a single contact
graph. We now partition the evidence into two independent graphs whose
agreement is the identification.

\begin{definition}[Instrument graph]
\label{def:inst-graph}
The \emph{instrument graph} $\Gph_I$ commits edges from:
\begin{itemize}[itemsep=1pt]
  \item Enzyme specificity (Catalyst~2)
  \item Fragment ion matching (Catalyst~3)
  \item Ladder topology (Catalyst~5)
  \item Fixed modification mass shifts (from sample preparation)
  \item Precursor isolation window and isotope envelope
\end{itemize}
These are \emph{preparation and acquisition facts}: they derive from
what was done to the sample and what the instrument measured.
\end{definition}

\begin{definition}[Mapping graph]
\label{def:map-graph}
The \emph{mapping graph} $\Gph_M$ commits edges from:
\begin{itemize}[itemsep=1pt]
  \item Precursor mass match against the database (Catalyst~1)
  \item S-Entropy coordinate distance (Catalyst~4)
  \item Mapping-set intersections (which proteins contain each
  candidate)
\end{itemize}
These are \emph{database and coordinate facts}: they derive from the
sequence database and the physicochemical property space.
\end{definition}

\begin{proposition}[Independence]
\label{prop:indep}
$\Gph_I$ and $\Gph_M$ share no mechanism: no mapping set encodes a
fragment match, and no fragment match encodes a homology relation.
\end{proposition}

\begin{construction}[Two-graph identification]
\label{con:twograph}
Compute $\amb_I(\Spec)$ and $\amb_M(\Spec)$ independently. The
identification is:
\[
  \amb(\Spec) \;=\; \amb_I(\Spec)\cap\amb_M(\Spec).
\]
By \Cref{thm:agreement}, this intersection is a relabelling invariant
of neither graph alone.
\end{construction}

\begin{proposition}[Disagreement is diagnostic]
\label{prop:disagree}
If $\amb_I\cap\amb_M=\emptyset$, inspect which is empty at $v$:
\begin{itemize}[itemsep=1pt]
  \item $\amb_M=\emptyset$: the candidate is absent from the database
  (novel peptide, contaminant, unexpected cleavage).
  \item $\amb_I=\emptyset$: the preparation chain is inconsistent
  with the candidate (wrong enzyme, missing modification).
\end{itemize}
\end{proposition}

% =====================================================================
\section{Isobaric Residues}
\label{sec:isobaric}
% =====================================================================

\begin{proposition}[Isobaric pairs are terminal regions]
\label{prop:isobaric}
Let $a,b$ be residues with equal mass (e.g.\ L and I, both
$113.084$~Da). If all committed edges are mass-derived, then
$\amb=\{a,b\}$ and closure holds. A catalyst narrows $\amb$ if and only
if it commits an edge whose weight differs between $a$ and $b$.
\end{proposition}

\begin{proof}
Mass-derived edges assign equal weights to $a$ and $b$, so their cut
keys agree. No further mass catalyst separates them. A non-mass catalyst
(ion mobility, high-energy CID producing $d$/$w$ ions,
EThcD~\citep{xiao2016distinguishing,zhokhov2017distinguishing})
commits an edge with a weight that differs, moving the claim region.
\end{proof}

\begin{remark}
The verdict \textsc{isobaric} is a \emph{correct closed answer}, not a
failure. It reports that the evidence committed cannot distinguish $a$
from $b$, and it states the condition under which further evidence would
narrow the region. This is more informative than reporting one of the
two with a flag.
\end{remark}

% =====================================================================
\section{Complexity}
\label{sec:complexity}
% =====================================================================

\begin{proposition}[Time complexity]
\label{prop:time}
The algorithm runs in
$O(|\Pep_M|\cdot(L\cdot m + |\Pep_M|\cdot|V|\cdot|E|))$ per spectrum,
where $|\Pep_M|$ is the precursor-filtered candidate count, $L$ is the
maximal peptide length, $m$ is the number of observed peaks,
and $|V|\cdot|E|$ is the max-flow cost per cut computation.
\end{proposition}

\begin{proof}
Precursor filtering: $O(\log|\Pep|+|\Pep_M|)$.
Fragment matching per candidate: $O(L\cdot m)$ if peaks are sorted
(binary search per theoretical ion) or $O(L+m)$ if both are sorted
(merge scan).
S-Entropy and topology: $O(L)$ per candidate.
Inter-item edges: $O(|\Pep_M|^2)$.
Max-flow per item: $O(|V|\cdot|E|)$ by Edmonds--Karp, where
$|V|=|\Pep_M|+1$ and $|E|\le|\Pep_M|+\binom{|\Pep_M|}{2}$.
Total: dominated by the max-flow computations over $|\Pep_M|$ items.
\end{proof}

\begin{remark}[Practical regime]
In the typical proteomics regime, $|\Pep_M|$ is $10$--$10^3$
(after precursor filtering at $10$~ppm on a human proteome),
$L\le 30$, $m\le 500$. The per-spectrum cost is dominated by
fragment matching at $O(|\Pep_M|\cdot L\cdot m)$, which is the
same asymptotic cost as conventional engines. The max-flow
computations on $|\Pep_M|+1$ vertices are fast for small candidate
sets and become the bottleneck only when $|\Pep_M|$ is large.
For $|\Pep_M|>10^3$, a pre-filter on fragment count (e.g.\ retain
only candidates with $n_{\mathrm{total}}\ge 3$) restores
tractability without compromising the closure guarantee, since
candidates with zero matched fragments are assigned inert catalytic
power and cannot enter a non-trivial claim region.
\end{remark}

% =====================================================================
\section{Properties of the Algorithm}
\label{sec:properties}
% =====================================================================

We now prove the structural properties claimed in the introduction.

\begin{theorem}[Termination]
\label{thm:termination}
The algorithm terminates in at most $D_0=\sum_{v\in\Pep_M}(|\amb(v)|-1)$
catalyst applications.
\end{theorem}

\begin{proof}
Each effective catalyst application strictly reduces $D_0$ by at least
$1$ (by the effective commitment assumption of the exclusion framework).
$D_0$ is a non-negative integer, so it reaches zero in at most $D_0$
steps. At that point all claim regions are singletons (closure) or no
effective catalyst remains (decline).
\end{proof}

\begin{theorem}[Closure strength]
\label{thm:closure-strength}
Closure of $\amb(\Spec)$ does not imply any particular confidence score,
and no confidence score implies closure. Specifically:
\begin{enumerate}[label=(\roman*),itemsep=1pt]
  \item A closed region with $|\amb|=2$ has maximum achievable
  confidence $1/2$ under any symmetric confidence functional.
  \item A high-confidence single candidate may be non-closed if a
  catalyst not yet applied would move the region.
\end{enumerate}
\end{theorem}

\begin{proof}
(i) is the L/I case: two candidates with equal cut keys under mass
evidence, no catalyst separating them. Any symmetric functional assigns
at most $1/2$ to each. (ii) is the usual situation in proteomics: a
high-scoring PSM whose FDR is computed without considering whether ion
mobility or alternative fragmentation would change the assignment.
\end{proof}

\begin{proposition}[Composite power profile]
\label{prop:composite}
The composite catalytic power of the five-catalyst chain, with
individual powers $\pot_1,\dots,\pot_5$, is:
\[
  \pot_{\mathrm{composite}}
  \;=\;
  1-\prod_{i=1}^{5}(1-\pot_i).
\]
Each additional catalyst closes a strictly smaller fraction of the
remaining gap. The marginal return of the $k$-th catalyst is
$(1-\pot_k)\prod_{i<k}(1-\pot_i)$ of the initial gap.
\end{proposition}

\begin{proof}
Direct application of \Cref{thm:mult}.
\end{proof}

\begin{proposition}[Decline is informative]
\label{prop:decline-informative}
Each decline verdict carries a structured payload:
\begin{itemize}[itemsep=1pt]
  \item \textsc{declined-no-candidate}: no database peptide has the
  observed precursor mass. This may indicate a peptide absent from
  the database, a novel modification, or a contaminant.
  \item \textsc{declined-no-evidence}: candidates exist but no fragment
  ions matched. This indicates poor spectrum quality or a fragmentation
  mode inconsistent with the expected ion types.
  \item \textsc{declined-unclosed}: a partial identification exists but
  an available catalyst (e.g.\ MS$^3$, ion mobility) was not applied.
  The payload names the blocking catalyst.
\end{itemize}
\end{proposition}

% =====================================================================
\section{Comparison with Conventional Database Search}
\label{sec:comparison}
% =====================================================================

\begin{table}[t]
\centering
\small
\caption{Structural comparison between conventional database search and
exclusion search.}
\label{tab:compare}
\begin{tabular}{@{}p{0.20\textwidth}p{0.35\textwidth}p{0.35\textwidth}@{}}
\toprule
\textbf{Property} & \textbf{Conventional} & \textbf{Exclusion search} \\
\midrule
Output & Ranked list of PSMs, filtered by FDR & Claim region
$\amb(\Spec)$: a set of candidates \\[3pt]
Termination & Score exceeds threshold & Closure: no available catalyst
moves the region \\[3pt]
Refusal & Silence (spectrum absent from output) & Structured verdict
with diagnostic payload \\[3pt]
Isobaric pairs & Arbitrary choice or quality flag & Correct terminal
region $\{L,I\}$ \\[3pt]
Confidence & External: FDR from decoy database & Internal: closure is
a property of the committed graph \\[3pt]
Evidence model & Score: a single number summarising the match &
Contact graph: a weighted graph with computable cuts \\[3pt]
Multiple evidence & Implicit in the score function & Explicit
catalyst chain with multiplicative law \\[3pt]
Database role & The search space & One evidence source (mapping graph);
instrument graph is independent \\
\bottomrule
\end{tabular}
\end{table}

% =====================================================================
\section{Implementation Notes}
\label{sec:implementation}
% =====================================================================

\subsection{Floor selection}

The floor $\flo$ is a modelling parameter. It must be strictly positive
and should reflect the minimum cost of any comparison. A natural choice
is:
\[
  \flo \;=\; \frac{\epsilon_f}{M_{\max}},
\]
where $\epsilon_f$ is the fragment mass tolerance and $M_{\max}$ is the
maximum precursor mass in the database. This ties the floor to the
instrument's resolving power.

\subsection{Max-flow implementation}

For small candidate sets ($|\Pep_M|\le 100$), a push-relabel
implementation~\citep{goldberg1988new} suffices. For larger sets, the
Stoer--Wagner algorithm~\citep{stoerwagner1997} computes the global
minimum cut without specifying source and sink, which is useful when
the goal is to find the item with the smallest separation cost.

\subsection{Parallelism}

Spectra are independent: the contact graph for one spectrum shares no
state with another. The algorithm parallelises trivially across spectra,
with each worker maintaining its own graph and max-flow solver.

\subsection{Tolerance calibration}

$\epsilon_p$ and $\epsilon_f$ are instrument-dependent. For Orbitrap
data, typical values are $\epsilon_p=10$~ppm and $\epsilon_f=0.02$~Da.
For Q-TOF, $\epsilon_f=0.1$~Da. The closure guarantee holds for any
positive tolerance; tighter tolerances yield smaller $|\Pep_M|$ and
faster convergence.

\subsection{Variable modifications}

Each variable modification on a residue creates a separate candidate in
the digest. For $k$ modification types on $n$ modifiable residues, this
multiplies the candidate space by up to $(k+1)^n$. The standard
mitigation---limiting the number of simultaneous modifications per
peptide---applies here as in conventional engines. The exclusion
framework adds one observation: a modification that changes S-Entropy
coordinates contributes evidence through Catalyst~4, whereas a
modification invisible to S-Entropy (e.g.\ a mass shift not altering
hydrophobicity, volume, or charge) is correctly treated as inert.

% =====================================================================
\section{Limitations}
\label{sec:limits}
% =====================================================================

\paragraph{The algorithm is not faster than conventional search.}
The per-spectrum cost is comparable to Sequest or Mascot for typical
candidate counts. The max-flow computations add overhead that
conventional engines do not pay. The advantage is structural (what the
output means) rather than computational (how fast it arrives).

\paragraph{The floor is a modelling choice.}
$\flo$ is not measured from the data. Different choices yield different
residuals. The theorems are $\flo$-independent; the numerical results
are not.

\paragraph{Closure is relative to available catalysts.}
A determination closed under mass-only catalysts may reopen when ion
mobility data becomes available. This is a feature---it is how the
system says ``more evidence would help''---but it means closure is not
permanence.

\paragraph{No FDR estimate.}
The algorithm does not compute an FDR. Closure is a different kind of
guarantee: it says ``no further available evidence can change this
answer'', not ``this answer is correct with probability $1-\alpha$''.
The two are complementary, and a system using exclusion search may
compute FDR separately from a decoy analysis of the closed regions.

\paragraph{S-Entropy coordinates are approximate.}
The coordinate assignments of \Cref{tab:aa-coords} are derived from
bulk physicochemical properties. In the context of a folded protein or
within a collision cell, local environment may shift effective
hydrophobicity or charge. The framework tolerates this via the
tolerance in catalyst 4's edge-weight function, but does not correct
for it.

\paragraph{Synthetic validation only.}
This manuscript describes the algorithm; it does not report a
benchmarking study. Performance on standard datasets
(e.g.\ iPRG 2012~\citep{chalkley2014proteomics}, NIST spectral
libraries) is the subject of a separate
evaluation.

% =====================================================================
\section{Conclusion}
\label{sec:conclusion}
% =====================================================================

We have given a complete, implementable database search algorithm for
peptide tandem mass spectrometry whose output is a region, whose
termination is closure, and whose refusal is a verdict. The algorithm
threads five catalysts---precursor mass, enzyme specificity, fragment
ion matching, S-Entropy coordinate similarity, and ladder topology
matching---through a per-spectrum contact graph, computes the claim
region by minimum cut, and checks closure against all available
catalysts. It terminates in bounded time, returns one of five verdicts,
and never produces a score without a verdict or a verdict without a
payload.

Three properties distinguish it from conventional search.

First, an isobaric pair such as $\{L,I\}$ is a correct closed answer
rather than a failure to resolve. The algorithm states the condition
under which a further catalyst (ion mobility, EThcD) would narrow the
region, and reports the region as-is when that catalyst is unavailable.

Second, a spectrum whose evidence is insufficient is not silently absent
from the output. It receives a decline verdict naming which evidence
source failed, which converts an uninterpretable silence into a
diagnosis.

Third, the independence of the instrument graph and the mapping graph
means the algorithm can distinguish two situations currently conflated
as ``unidentified'': the peptide is absent from the database (mapping
graph empty) versus the preparation chain could not have produced the
candidate (instrument graph empty).

What the algorithm does not provide is a probability. Closure and FDR
are orthogonal guarantees, and a complete system may well compute both.
The contribution here is the structural layer---what the search returns
and when it stops---on which any scoring function or statistical model
may be placed.

% =====================================================================
% References
% =====================================================================
\begingroup
\small
\begin{thebibliography}{99}

\bibitem{eng1994approach}
J.~K. Eng, A.~L. McCormack, and J.~R. Yates~III.
\newblock An approach to correlate tandem mass spectral data of peptides
  with amino acid sequences in a protein database.
\newblock \emph{J.\ Am.\ Soc.\ Mass Spectrom.}, 5(11):976--989, 1994.

\bibitem{perkins1999probability}
D.~N. Perkins, D.~J.~C. Pappin, D.~M. Creasy, and J.~S. Cottrell.
\newblock Probability-based protein identification by searching sequence
  databases using mass spectrometry data.
\newblock \emph{Electrophoresis}, 20(18):3551--3567, 1999.

\bibitem{cox2008maxquant}
J.~Cox and M.~Mann.
\newblock {MaxQuant} enables high peptide identification rates,
  individualized p.p.b.-range mass accuracies and proteome-wide protein
  quantification.
\newblock \emph{Nat.\ Biotechnol.}, 26(12):1367--1372, 2008.

\bibitem{elias2007target}
J.~E. Elias and S.~P. Gygi.
\newblock Target-decoy search strategy for increased confidence in
  large-scale protein identifications by mass spectrometry.
\newblock \emph{Nat.\ Methods}, 4(3):207--214, 2007.

\bibitem{kall2007semi}
L.~K\"all, J.~D. Storey, M.~J. MacCoss, and W.~S. Noble.
\newblock Semi-supervised learning for peptide identification from
  shotgun proteomics datasets.
\newblock \emph{Nat.\ Methods}, 4(11):923--925, 2007.

\bibitem{the2016fast}
M.~The, M.~J. MacCoss, W.~S. Noble, and L.~K\"all.
\newblock Fast and accurate protein false discovery rates on large-scale
  proteomics data sets with {Percolator 3.0}.
\newblock \emph{J.\ Am.\ Soc.\ Mass Spectrom.}, 27(11):1719--1727, 2016.

\bibitem{paizs2005fragmentation}
B.~Paizs and S.~Suhai.
\newblock Fragmentation pathways of protonated peptides.
\newblock \emph{Mass Spectrom.\ Rev.}, 24(4):508--548, 2005.

\bibitem{olsen2004trypsin}
J.~V. Olsen, S.-E. Ong, and M.~Mann.
\newblock Trypsin cleaves exclusively {C}-terminal to arginine and
  lysine residues.
\newblock \emph{Mol.\ Cell.\ Proteomics}, 3(6):608--614, 2004.

\bibitem{krokhin2006sequence}
O.~V. Krokhin.
\newblock Sequence-specific retention calculator.\ {Algorithm} for
  peptide retention prediction in ion-pair {RP-HPLC}.
\newblock \emph{Anal.\ Chem.}, 78(22):7785--7795, 2006.

\bibitem{kyte1982simple}
J.~Kyte and R.~F. Doolittle.
\newblock A simple method for displaying the hydropathic character of a
  protein.
\newblock \emph{J.\ Mol.\ Biol.}, 157(1):105--132, 1982.

\bibitem{zamyatnin1972protein}
A.~A. Zamyatnin.
\newblock Protein volume in solution.
\newblock \emph{Prog.\ Biophys.\ Mol.\ Biol.}, 24:107--123, 1972.

\bibitem{fordfulkerson1956}
L.~R. Ford and D.~R. Fulkerson.
\newblock Maximal flow through a network.
\newblock \emph{Can.\ J.\ Math.}, 8:399--404, 1956.

\bibitem{goldberg1988new}
A.~V. Goldberg and R.~E. Tarjan.
\newblock A new approach to the maximum-flow problem.
\newblock \emph{J.\ ACM}, 35(4):921--940, 1988.

\bibitem{stoerwagner1997}
M.~Stoer and F.~Wagner.
\newblock A simple min-cut algorithm.
\newblock \emph{J.\ ACM}, 44(4):585--591, 1997.

\bibitem{xiao2016distinguishing}
Y.~Xiao, M.~M. Vecchi, and D.~Wen.
\newblock Distinguishing between leucine and isoleucine by integrated
  {LC-MS} analysis using an {Orbitrap Fusion} mass spectrometer.
\newblock \emph{Anal.\ Chem.}, 88(21):10757--10766, 2016.

\bibitem{zhokhov2017distinguishing}
S.~S. Zhokhov, S.~V. Kovalyov, A.~S. Samgina, and A.~T. Lebedev.
\newblock An {EThcD}-based method for discrimination of leucine and
  isoleucine residues in tryptic peptides.
\newblock \emph{J.\ Am.\ Soc.\ Mass Spectrom.}, 28(8):1600--1611, 2017.

\bibitem{tran2017novo}
N.~H. Tran, X.~Zhang, L.~Xin, B.~Shan, and M.~Li.
\newblock De novo peptide sequencing by deep learning.
\newblock \emph{Proc.\ Natl.\ Acad.\ Sci.\ USA}, 114(31):8247--8252,
  2017.

\bibitem{zhou2017pdeep}
X.-X. Zhou, W.-F. Zeng, H.~Chi, C.~Luo, C.~Liu, \emph{et al.}
\newblock {pDeep}: predicting {MS/MS} spectra of peptides with deep
  learning.
\newblock \emph{Anal.\ Chem.}, 89(23):12690--12697, 2017.

\bibitem{chalkley2014proteomics}
R.~J. Chalkley, S.~P. Huang, M.~D. Guttman, and S.~Koyama.
\newblock Comprehensive analysis of a multidimensional liquid
  chromatography mass spectrometry dataset acquired on a quadrupole
  selecting, quadrupole collision cell, time-of-flight mass spectrometer:
  {II}.\ New developments in {Protein Prospector} allow for reliable and
  comprehensive automatic analysis of large datasets.
\newblock \emph{Mol.\ Cell.\ Proteomics}, 4(8):1194--1204, 2005.

\bibitem{sachikonye2025peptide}
K.~F. Sachikonye.
\newblock Peptide identity as the termination of exclusion: mass
  invariance, promiscuous intersection, and a stopping criterion for
  sequence determination.
\newblock Manuscript, 2025.

\bibitem{sachikonye2025molecular}
K.~F. Sachikonye.
\newblock On the consequences of categorical completion on biological
  sequences: thermodynamic symbolic computational molecular language for
  database-free peptide sequence reconstruction.
\newblock Manuscript, 2025.

\bibitem{sachikonye2025fragmentation}
K.~F. Sachikonye.
\newblock Categorical fragmentation networks in tandem mass spectrometry:
  phase-lock topology and entropy-intensity relations in peptide
  fragmentation.
\newblock Manuscript, 2025.

\bibitem{roepstorff1984proposal}
P.~Roepstorff and J.~Fohlman.
\newblock Proposal for a common nomenclature for sequence ions in mass
  spectra of peptides.
\newblock \emph{Biomed.\ Mass Spectrom.}, 11(11):601, 1984.

\bibitem{dancik1999novo}
V.~Dan\v{c}\'ik, T.~A. Addona, K.~R. Clauser, J.~E. Vath, and
  P.~A. Pevzner.
\newblock De novo peptide sequencing via tandem mass spectrometry.
\newblock \emph{J.\ Comput.\ Biol.}, 6(3--4):327--342, 1999.

\bibitem{ma2003peaks}
B.~Ma, K.~Zhang, C.~Hendrie, C.~Liang, M.~Li, A.~Doherty-Kirby,
  and G.~Lajoie.
\newblock {PEAKS}: powerful software for peptide de novo sequencing by
  tandem mass spectrometry.
\newblock \emph{Rapid Commun.\ Mass Spectrom.}, 17(20):2337--2342, 2003.

\bibitem{fenn1989electrospray}
J.~B. Fenn, M.~Mann, C.~K. Meng, S.~F. Wong, and C.~M. Whitehouse.
\newblock Electrospray ionization for mass spectrometry of large
  biomolecules.
\newblock \emph{Science}, 246(4926):64--71, 1989.

\end{thebibliography}
\endgroup

\end{document}